\section{Spatial moments and reduction to a coordinate direction}
\label{sec:radius}

The finite mean width from Proposition~\ref{prop:mean-width} controls a
regeneration word only in its direction of progress.  We now show that
coexistence of the two escape signs also controls its spatial extent.
This will give two limiting rays and reduce the problem to a coordinate
direction.  The radius argument adapts
\cite[Proposition~3.1]{OpenAI2026}, with a marked connector and a
separated infinite continuation supplying the finite-range version.

For a regeneration word $\gamma$ with position sequence
$(x_0,\ldots,x_m)$, written relative to $x_0=0$, define its radius by
\[
  \mathcal R(\gamma)=\max_{0\le j\le m}|x_j|,
\]
where $|\cdot|$ denotes Euclidean norm.  In particular, the terminal
arrival contributes to the radius.  Distances used for finite-range
independence remain $\ell^\infty$ distances.

\subsection{An integrable spatial radius}

We first record an elementary estimate that controls all initial sums at
once.  Stationarity, rather than independence, is enough for this estimate.
It is the stationary-sequence form of Hopf's maximal ergodic inequality;
see Garsia~\cite{Garsia1965}.  We include a disjoint-interval proof.

\begin{lemma}[A maximal estimate for initial averages]
\label{lem:stationary-maximal}
Let $(Z_i)_{i\ge1}$ be a stationary sequence of nonnegative integrable
random variables.  For every integer $n\ge1$ and every $a>0$,
\[
  \PP\left\{\max_{1\le k\le n}\frac1k\sum_{i=1}^k Z_i>a\right\}
  \le \frac{\EE Z_1}{a}.
\]
\end{lemma}

\begin{proof}
Call an integer $j\ge1$ bad if some interval starting at $j$, of length
at most $n$, has average greater than $a$.  Among the bad starts in
$\{1,\ldots,m\}$, select the leftmost one and a witnessing interval.
Continue by selecting the leftmost bad start strictly to the right of
the selected interval, and repeat.  These disjoint intervals cover all
bad starts in $\{1,\ldots,m\}$ and lie in $\{1,\ldots,m+n\}$.
Their total length is at least the number of those bad starts, while
the sum of the $Z_i$ over them is at least $a$ times their total length.
Consequently
\[
  a\,\#\{j\le m:j\text{ is bad}\}\le\sum_{i=1}^{m+n}Z_i.
\]
Taking expectations, using stationarity, dividing by $m$, and letting
$m\to\infty$ proves the assertion.
\end{proof}

\begin{proposition}[Radius moment under coexistence]
\label{prop:radius-moment}
Let $v\ne0$ satisfy $\|v\|_\infty=1$.  If
$P_0(A_v)>0$ and $P_0(A_{-v})>0$, then the regeneration-word laws of
Proposition~\ref{prop:cut-law} satisfy
\begin{equation}
  E_{\nu_v}\mathcal R+E_{\nu_{-v}}\mathcal R<\infty.
  \label{eq:radius-moment}
\end{equation}
\end{proposition}

\begin{proof}
Both nonbacktracking probabilities $p_v,p_{-v}$ are positive by
Proposition~\ref{prop:sign-union}.  For $\sigma\in\{+,-\}$ write
\[
  P_\sigma^*=\widehat P_{\sigma v},\qquad
  E_\sigma^*=E_{P_\sigma^*}.
\]
Here $\sigma v$ means $v$ or $-v$, respectively.  Under either law, let
$L_i,\mathcal R_i$ be the width and radius of its $i$th regeneration
word, and put
\[
  H_0=0,\qquad H_k=\sum_{i=1}^k L_i.
\]
An unsubscripted $L$ or $\mathcal R$ denotes the corresponding variable
for one word.  All constants chosen below are independent of the scale
parameters $i_0$ and $a$.

Suppose that \eqref{eq:radius-moment} fails.  The truncated mean
\[
  I(t)=\sum_{\sigma\in\{+,-\}}E_\sigma^*\min(\mathcal R,t),
  \qquad t\ge0,
\]
is then nondecreasing and unbounded.  It is concave, $I(0)=0$, and
$I(t)=o(t)$: the last assertion follows by dominated convergence from
the almost-sure finiteness of each regeneration word.

We will obtain a fixed positive probability that the final supremum of
one of the two signed heights lies in a bounded window, and then place
these windows arbitrarily far apart.  The unbounded truncated mean
will produce a large spatial excursion after a controlled initial
segment, and hence a record for a tilted linear functional.  A marked
connector will let us append a separated continuation in the opposite
direction while preserving the height maximum already attained.

\paragraph{The scale and a controlled initial segment.}
Proposition~\ref{prop:mean-width} and the strong law allow us to choose
$c,C>0$ and $C_0\ge0$ so that, under each sign law, with probability
greater than $0.95$,
\begin{equation}
  ck-C_0\le H_k\le Ck+C_0\qquad(k\ge0).
  \label{eq:radius-height-bounds}
\end{equation}
Indeed, choose $c$ below both mean widths and $C$ above both, and then
choose a deterministic $C_0$ controlling the finitely many initial
deviations with the stated probability.  Next choose an integer
$C_1\ge1$, a number $b>1/c$, a number $B>1$, and $\eta>0$, in that
order, with
\begin{equation}
  C_1c>C+4,\qquad
  \frac{B-1}{\sqrt d}-b(C+2)>2,\qquad
  4C_1\eta<0.15.
  \label{eq:radius-constants}
\end{equation}

For $a>0$ sufficiently large, let $N=N(a)$ be the first positive integer
such that
\[
  I(aN)\ge\eta a.
\]
It exists because $I$ is unbounded, and $N(a)\to\infty$ because
$I(t)=o(t)$.  For large $a$, minimality and concavity give
\begin{equation}
  \eta a\le I(aN)
  \le\frac{N}{N-1}I(a(N-1))<2\eta a.
  \label{eq:radius-scale}
\end{equation}

Let $\mathcal J_m$ be the event that
\eqref{eq:radius-height-bounds} holds through $m$ and
\[
  \sum_{i=1}^k\mathcal R_i\le ak\qquad(1\le k\le m).
\]
For both signs and all sufficiently large $a$,
\begin{equation}
  P_\sigma^*(\mathcal J_{C_1N})>0.8.
  \label{eq:radius-controlled}
\end{equation}
To see this, set $n=C_1N$ and truncate each radius at $an$.
Its mean is at most
\[
  I(an)\le C_1 I(aN)<2C_1\eta a.
\]
Lemma~\ref{lem:stationary-maximal} bounds by $2C_1\eta$ the probability
that any initial average of the truncated radii, through $n$, exceeds
$a$.  Also
\[
  P_\sigma^*\{\mathcal R_i>an\text{ for some }i\le n\}
  \le nP_\sigma^*(\mathcal R>an)
  \le\frac{E_\sigma^*\min(\mathcal R,an)}a
  <2C_1\eta.
\]
Combining these estimates with the probability of
\eqref{eq:radius-height-bounds} proves
\eqref{eq:radius-controlled}.

\paragraph{A large excursion following that segment.}
Choose an integer $i_0$ large enough that
\begin{equation}
\begin{gathered}
  \sum_\sigma\sum_{i\ge i_0}P_\sigma^*(L>i)<\frac{\eta}{4B},\\
  c(i-1)-C_0\ge\frac{ci}{2},\qquad
  C(i-1)+C_0+i\le(C+2)i\qquad(i\ge i_0).
\end{gathered}
\label{eq:radius-initial-index}
\end{equation}
The first condition follows from integrability of the widths.
With $i_0$ fixed, take $a$ sufficiently large that $N\ge i_0$ and
$I(Bai_0)<\eta a/4$, as well as all preceding large-$a$ requirements.
For $i_0\le i\le N$ define
\[
  \mathcal L_i=\{\mathcal R_i>Bai,\ L_i\le i\},\qquad
  \Lambda_\sigma=\sum_{i=i_0}^N P_\sigma^*(\mathcal L_i).
\]
Integration of the decreasing radius tails gives
\begin{align*}
  \sum_\sigma\sum_{i=i_0}^N P_\sigma^*(\mathcal R>Bai)
  &\ge\frac{I(Ba(N+1))-I(Bai_0)}{Ba},\\
  \sum_\sigma\sum_{i=1}^N P_\sigma^*(\mathcal R>Bai)
  &\le\frac{I(BaN)}{Ba}.
\end{align*}
The first right-hand side is at least $3\eta/(4B)$, by
\eqref{eq:radius-scale} and monotonicity of $I$.  The second is at most
$2\eta$, by concavity.  Subtracting the width-tail bound in
\eqref{eq:radius-initial-index} therefore yields
\begin{equation}
  \frac{\eta}{2B}\le\Lambda_++\Lambda_-\le2\eta.
  \label{eq:radius-large-excursion-mass}
\end{equation}

Choose a sign $\sigma$ for which $\Lambda_\sigma\ge\eta/(4B)$, and count
the indices for which both $\mathcal J_{i-1}$ and $\mathcal L_i$ occur.
Independence of the $i$th word from its predecessors and
\eqref{eq:radius-controlled} give a mean of at least
$0.8\Lambda_\sigma$.  The second moment is at most
$\Lambda_\sigma+\Lambda_\sigma^2$, by domination by the sum of the
independent indicators $\one_{\mathcal L_i}$.
Cauchy--Schwarz and \eqref{eq:radius-large-excursion-mass} imply that,
with probability at least a constant $\delta>0$ independent of $i_0,a$,
there is an index $i\in[i_0,N]$ with
\begin{equation}
  \mathcal J_{i-1}\cap\mathcal L_i.
  \label{eq:radius-selected-excursion}
\end{equation}

\paragraph{A record in a tilted direction.}
For this sign write $h(x)=\sigma v\cdot x$.  On
\eqref{eq:radius-selected-excursion}, all preceding positions have norm
at most $a(i-1)$, whereas some position $X$ in word $i$ satisfies
$|X|>(B-1)ai$.  Hence, for some coordinate step $u\in U$,
\[
  u\cdot X>\frac{(B-1)ai}{\sqrt d}.
\]
Throughout this word, the running maximum of $h$ lies between
$H_{i-1}$ and $H_i$, and therefore in
$[ci/2,(C+2)i]$.  Define the linear functional
\[
  Y(x)=u\cdot x-ba h(x).
\]
Before word $i$, nonnegativity of $h$ gives $Y\le a(i-1)$.
At the indicated position, \eqref{eq:radius-constants} gives $Y>2ai$.
There is consequently a strict $Y$-record whose running $h$-maximum
belongs to the deterministic window
\begin{equation}
  \mathcal W=[ci_0/2,(C+2)N].
  \label{eq:radius-window}
\end{equation}

For each fixed $u$, let $T$ be the first positive time at which the
position is a strict $Y$-record, the running $h$-maximum belongs to
$\mathcal W$, and the entire prefix has nonnegative $h$-heights.
These conditions make $T$ a stopping time.  Since there are $2d$
choices of $u$, one of them satisfies
$P_\sigma^*(T<\infty)\ge\delta/(2d)$.  Fix that choice and set
$p_* = \min(p_v,p_{-v})$.  Under the raw law,
\begin{equation}
  P_0(T<\infty)\ge\frac{p_*\delta}{2d}.
  \label{eq:radius-record-probability}
\end{equation}

We have found a record with a uniformly positive probability, while its
running height maximum lies in a window that can be moved arbitrarily
far out.  We next continue from this record so that the final height
maximum stays in the same window.  The continuation must remain
separated from the observed prefix for its annealed probability to
factor.

\paragraph{A separated continuation in the opposite direction.}
Choose a coordinate step $e$ with $h(e)=-1$.  From $X_T$ force $m_0$
successive steps $e$, all with mark $0$, where $m_0$ is a fixed positive
integer satisfying
\begin{equation}
  m_0b>1+bC_0+Rdb+1.
  \label{eq:radius-connector-length}
\end{equation}
The choice of $m_0$ depends only on the fixed constants $b,C_0,R,d$,
and is independent of $i_0,a$.
These steps have weight $\lambda^{m_0}$.  At their endpoint
$z=X_T+m_0e$, consider the translated raw event consisting of
$D_{-\sigma v}$ and $\mathcal J_{C_1N}$ for the opposite regeneration
sequence, with the required cuts finite.  Its raw probability is at
least $0.8p_*$ by \eqref{eq:radius-controlled} and
Proposition~\ref{prop:cut-law}.

Every position of the continuation from $z$ is separated by more than
$R$ from the prefix through $T$.  There are two parts to this geometric
assertion.
During opposite word $k+1$, for $0\le k<C_1N$, displacement from $z$
in the $u$ coordinate is at least $-a(k+1)$, by the initial-sum radius
bound in $\mathcal J_{C_1N}$; displacement in $h$ is at most $-ck+C_0$.
Its total $Y$-increment from $X_T$, including
the forced steps, is therefore at least
\begin{align}
  &m_0(ba-1)-a(k+1)+ba(ck-C_0)\notag\\
  &\hspace{12mm}
  =a\bigl(m_0b-1-bC_0+(bc-1)k\bigr)-m_0.
  \label{eq:radius-tilted-gap}
\end{align}
Since $bc>1$, \eqref{eq:radius-connector-length} makes this quantity
greater than $R(1+dba)$ for all sufficiently large $a$, uniformly in
$k$.  Every prefix position has $Y$-value at most $Y(X_T)$, while
\[
  |Y(x)-Y(y)|\le(1+dba)\|x-y\|_\infty.
\]
Thus the first $C_1N$ opposite words are at distance greater than $R$
from every prefix position.

At the end of these words the absolute $h$-coordinate is at most
\begin{equation}
  (C+2)N-m_0-cC_1N+C_0<-dR
  \label{eq:radius-tail-gap}
\end{equation}
for large $a$, by \eqref{eq:radius-constants}.  All subsequent positions
stay at or below this height, since this endpoint is a cut in direction
$-\sigma v$.  The prefix has nonnegative $h$-heights and
$|h(x)-h(y)|\le d\|x-y\|_\infty$, so the entire remaining tail is
also separated from the prefix.  The tilted inequality protects the
initial opposite words; the height inequality protects everything
after their last cut.  No bound on $Y$ is required for this remaining tail.

We now justify the probability calculation for this infinite
continuation.  Fix a finite marked prefix $\gamma$ realizing $T$, and
let $A_\gamma$ be its set of positions, including its terminal arrival.
The preceding deterministic bounds show that every path in the
specified continuation event avoids the closed $R$-neighborhood of
$A_\gamma$.  Consequently its quenched probability can be computed
using the walk killed on first arrival in that neighborhood.  By
Lemma~\ref{lem:separated-weights}, this probability depends only on
rows outside the neighborhood, even though the event specifies future
cuts and an infinite nonbacktracking tail.  Those rows are separated
from the active departures of $\gamma$.

The forced connector has constant quenched weight and reveals no rows.
The weight of $\gamma$, the connector, and the continuation therefore
factors as
\[
  W(\gamma)\lambda^{m_0}
  P_0\{D_{-\sigma v},\ \mathcal J_{C_1N}
                 \text{ for its regeneration sequence}\}
  \ge 0.8p_*\lambda^{m_0}W(\gamma).
\]
Here stationarity identifies the last factor with the raw probability
from $z$, and the almost-sure existence of the indicated cuts under
$D_{-\sigma v}$ is understood.  This use of separated rows is made
after fixing $\gamma$; no independence for a randomly chosen region is
assumed.

The connector decreases $h$, and the opposite continuation stays at
or below its own initial height.  Thus the final supremum of $h$ is
exactly the running supremum at $T$ and belongs to $\mathcal W$.
Summing over the disjoint prefixes realizing $T$ and using
\eqref{eq:radius-record-probability} yields
\begin{equation}
  P_0\left\{\sup_{n\ge0}\sigma v\cdot X_n\in\mathcal W\right\}
  \ge\rho,
  \qquad
  \rho:=\frac{0.8\lambda^{m_0}p_*^2\delta}{2d}>0.
  \label{eq:radius-final-supremum}
\end{equation}

Finally, choose $i_0$ and then $a$ successively so that the resulting
windows \eqref{eq:radius-window} are pairwise disjoint.  This is
possible by taking each new lower endpoint beyond the preceding upper
endpoint before choosing its $a$.  The sign $\sigma$ may change from
one window to the next, but for either fixed sign the corresponding
final-supremum events are disjoint.  Their probabilities, summed over
all windows and both signs, are at most $2$, contradicting the uniform
lower bound \eqref{eq:radius-final-supremum}.  This proves
\eqref{eq:radius-moment}.
\end{proof}

\subsection{Opposite limiting rays}

The radius moment controls the path between cuts, so the ordinary
strong law for regeneration displacements also describes the full
trajectory.  The limiting-ray argument and the coordinate reduction
follow \cite[Section~4.1]{OpenAI2026}.

\begin{lemma}[Limiting rays under coexistence]
\label{lem:opposite-rays}
Under the coexistence hypotheses of Proposition~\ref{prop:radius-moment},
there are deterministic unit vectors $r_+,r_-$ such that, for
$\sigma\in\{+,-\}$,
\[
  |X_n|\longrightarrow\infty,\qquad
  \frac{X_n}{|X_n|}\longrightarrow r_\sigma
  \quad P_0\text{-almost surely on }A_{\sigma v}.
\]
They satisfy $r_\sigma\cdot(\sigma v)>0$ and $r_-=-r_+$.
\end{lemma}

\begin{proof}
The terminal displacement of a word of law $\nu_{\sigma v}$ is
integrable by Proposition~\ref{prop:radius-moment}.  Write its mean as
$b_\sigma'\in\RR^d$.  Since
\[
  b_\sigma'\cdot(\sigma v)=E_{\nu_{\sigma v}}L>0,
\]
this vector is nonzero.  Under $\widehat P_{\sigma v}$, the $k$th cut
location divided by $k$ tends almost surely to $b_\sigma'$.
Integrability also gives $\mathcal R_k/k\to0$ almost surely: for each
$\varepsilon>0$, the sum of
$\widehat P_{\sigma v}(\mathcal R_k>\varepsilon k)$ is finite, and
Borel--Cantelli applies.  Between consecutive cuts the displacement
from the earlier cut is bounded by the next radius.  Since the number
of completed words tends to infinity along the trajectory, it follows
that $|X_n|\to\infty$ and
\[
  \frac{X_n}{|X_n|}\longrightarrow
  r_\sigma:=\frac{b_\sigma'}{|b_\sigma'|}.
\]
Under the raw law on $A_{\sigma v}$, the first cut is finite and its
translated suffix has law $\widehat P_{\sigma v}$ by
Proposition~\ref{prop:cut-law}; hence the same conclusion holds there.

The signs of the projections on $v$ make $r_+$ and $r_-$ distinct.
Suppose they are not opposite.  Then $w=r_++r_-$ is nonzero and
\[
  w\cdot r_+=w\cdot r_-=1+r_+\cdot r_->0.
\]
Proposition~\ref{prop:sign-union} gives
$P_0(A_v\cup A_{-v})=1$, so the limiting rays imply $P_0(A_w)=1$.
Rescale $w$ to have $\ell^\infty$ norm one and apply the cut
construction in that direction.

Under $\widehat P_w$, the limiting ray exists almost surely because
this law is a conditioning of the raw law just considered.  Its value
is unchanged on deleting any finite number of regeneration words and
translating the remaining path: such operations do not affect a
limiting direction when positions tend to infinity in norm.
Consequently it is a tail random variable of the iid words under
$\widehat P_w$: for every $k$, its value can be computed from the words
after $k$.  The independent tail zero--one law makes it
deterministic.  Since $P_0(A_w)=1$, the raw walk has a first cut in
direction $w$ almost surely, with translated suffix law
$\widehat P_w$.  Its limiting ray must therefore have that same
deterministic value almost surely.  This contradicts the two distinct
rays $r_+,r_-$, each occurring with positive raw probability.  Hence
$r_-=-r_+$.
\end{proof}

\begin{proposition}[Reduction to a coordinate direction]
\label{prop:coordinate-reduction}
If $0<P_0(A_v)<1$ for some fixed nonzero real direction $v$, then there
is a coordinate $i\in\{1,\ldots,d\}$ such that
\[
  P_0(A_{e_i})>0\qquad\text{and}\qquad P_0(A_{-e_i})>0.
\]
\end{proposition}

\begin{proof}
After normalizing $v$, Proposition~\ref{prop:sign-union} shows that
both signs of escape in direction $v$ have positive probability.
Choose a coordinate with nonzero projection on the opposite rays
from Lemma~\ref{lem:opposite-rays}.  Along one ray this coordinate
tends to $+\infty$, and along the other it tends to $-\infty$.
Each ray occurs with positive probability, giving the assertion.
\end{proof}

To prove Theorem~\ref{thm:main}, it therefore remains to exclude
coexistence in a coordinate direction.  After permuting coordinates,
we may take that direction to be $e_1$.
